\documentclass[11pt,aspectratio=169]{beamer}

\usetheme{Madrid}
\usecolortheme{seahorse}

\usepackage[utf8]{inputenc}
\usepackage{amsmath,amssymb}
\usepackage{algpseudocode}
\usepackage{algorithmicx}
\usepackage{xcolor}
\usepackage{listings}
\usepackage{booktabs}
\usepackage{tikz}
\usetikzlibrary{arrows.meta,positioning,fit,backgrounds}

\setbeamertemplate{navigation symbols}{}
\setbeamertemplate{footline}[frame number]

\lstdefinestyle{py}{
  language=Python,
  basicstyle=\ttfamily\scriptsize,
  keywordstyle=\color{blue!70!black}\bfseries,
  commentstyle=\color{gray},
  stringstyle=\color{green!40!black},
  showstringspaces=false,
  breaklines=true,
  numbers=left,
  numberstyle=\tiny\color{gray},
  frame=single,
  xleftmargin=1.5em,
  aboveskip=0.3em,
  belowskip=0.3em
}

\title{Beyond the Syllabus}
\subtitle{Greedy Correctness, Matroids, and Online Algorithms with Predictions\\
\small Week 4 Supplement -- TOAE209/TOAH209: Design and Analysis of Algorithms}
\author{Foreign Trade University}
\date{}

\begin{document}

\begin{frame}
  \titlepage
\end{frame}

\begin{frame}{Outline}
  \tableofcontents
\end{frame}

% =================================================================
\section{Motivation}
% =================================================================

\begin{frame}{Why look beyond the syllabus?}
  \begin{itemize}
    \item This week you proved earliest-finish-time-first optimal for interval
    scheduling using an \textbf{exchange argument}, and saw the same proof style reused
    for interval partitioning, minimizing lateness, and job sequencing.
    \item Three questions a curious student should ask next:
    \begin{enumerate}
      \item Exchange arguments feel ad hoc, tailored to each problem. Is there a single
      \textbf{structural condition} on the underlying combinatorial problem that tells
      you, in advance, whether \emph{any} greedy rule can possibly be optimal?
      \item Every algorithm this week assumed you see the \emph{whole} instance before
      sorting. What happens when requests arrive \textbf{online}, one at a time, and you
      must accept or reject immediately?
      \item If you had a (possibly imperfect) \textbf{prediction} about future requests
      -- say, from a machine-learning forecast -- could you beat the classical online
      lower bounds?
    \end{enumerate}
  \end{itemize}
  \begin{alertblock}{How this supplement is different from a survey}
    Every citation below was checked against its actual paper/venue before this slide
    was written. Where a specific bound or year is quoted, it is sourced to a specific,
    verifiable paper.
  \end{alertblock}
\end{frame}

% =================================================================
\section{Generalizing Greedy Correctness: Matroids}
% =================================================================

\begin{frame}{When does \emph{any} greedy rule work?}
  \begin{block}{Rado--Edmonds theorem (classical)}
    Let $(E, \mathcal{I})$ be an independence system (a ground set $E$ with a family
    $\mathcal{I}$ of ``independent'' subsets, downward closed). For \textbf{every}
    choice of nonnegative weights, the greedy algorithm (sort by weight, add an element
    if it keeps the current set independent) finds a maximum-weight independent set
    \textbf{if and only if} $(E, \mathcal{I})$ is a \textbf{matroid}: $\mathcal{I}$ is
    hereditary, and whenever $A, B \in \mathcal{I}$ with $|A| < |B|$, some $x \in
    B\setminus A$ has $A \cup \{x\} \in \mathcal{I}$.
  \end{block}
  \begin{itemize}
    \item This is exactly why earliest-finish-time-first worked: the exchange step in
    this week's proof (swap a piece of $O$ for a piece of $G$ without losing
    feasibility) is a special case of the matroid exchange property.
    \item It also explains \emph{why} some greedy rules fail: the ``compatible
    intervals'' independence system is close to, but not literally, a matroid in
    general -- which is why the correctness proof needs the specific finish-time
    ordering, not an arbitrary one.
  \end{itemize}
\end{frame}

\begin{frame}{From matroids to submodular functions}
  \begin{itemize}
    \item Many modern applications (sensor placement, feature selection, influence
    maximization) replace ``maximize total weight'' with maximize a
    \textbf{submodular} set function $f$: adding an element helps \emph{less} as the
    set grows (diminishing returns), formally $f(A\cup\{x\}) - f(A) \ge f(B\cup\{x\}) -
    f(B)$ for $A \subseteq B$.
    \item For a single cardinality constraint ($|A|\le k$), the plain greedy
    (repeatedly add the element with the largest marginal gain) is only a
    $(1-1/e)$-approximation, not exact -- diminishing returns break the exchange
    argument's ``no worse off'' step.
    \item Under a general \textbf{matroid} constraint (not just $|A|\le k$), the
    continuous-greedy method of Cal\i nescu, Chekuri, P\'al \& Vondr\'ak (2011) still
    achieves $(1-1/e)$, matching the best possible unless $P=NP$.
  \end{itemize}
\end{frame}

\begin{frame}{A very recent refinement: how many oracle calls does it take?}
  \small
  \begin{block}{Kobayashi \& Terao, ICALP 2024 -- ``Subquadratic Submodular
  Maximization with a General Matroid Constraint''}
    Achieves a $(1-1/e-\varepsilon)$-approximation using only
    $\widetilde{O}_\varepsilon(\sqrt{r}\,n)$ independence-oracle and value-oracle
    queries, improving the previous best bound of
    $\widetilde{O}_\varepsilon(r^2 + \sqrt{r}\,n)$ due to Buchbinder, Feldman \&
    Schwartz. (LIPIcs vol.\ 297, pp.\ 100:1--100:19; arXiv:2405.00359.)
  \end{block}
  \begin{itemize}
    \item Here $n$ = ground-set size, $r$ = matroid rank. Concretely, at
    $n = 10^6, r = 10^5$: the old bound's $r^2$ term is $10^{10}$, while the new bound's
    $\sqrt{r}\,n \approx 3.16\times 10^8$ -- a $\sim 30\times$ reduction in the number
    of expensive feasibility/value checks, for the \emph{same} approximation guarantee.
    \item Same flavor as this week's lesson: the \emph{quality} bound ($1-1/e$) was
    already known; the 2024 advance is entirely about \textbf{how cheaply} you can
    reach it -- exactly the ``prove it correct, then make it fast'' arc this course
    follows.
  \end{itemize}
\end{frame}

% =================================================================
\section{Greedy Without Seeing the Future: Online Predictions}
% =================================================================

\begin{frame}{The online version of interval scheduling}
  \begin{itemize}
    \item This week's algorithm assumes all $n$ requests are known before you sort by
    finish time. In an \textbf{online} version, requests arrive one at a time and you
    must accept or reject \emph{immediately}, with no take-backs.
    \item A short adversary argument shows no deterministic online algorithm can
    achieve a constant competitive ratio in general: the adversary can always follow
    an accepted long interval with many short, mutually compatible intervals that the
    long one made infeasible to accept, and can do this with unboundedly bad ratio.
    (This is a standard, easy folklore argument in online algorithms -- unlike the
    dated results below, it needs no specific citation.)
    \item So classical worst-case online algorithms for this exact problem are
    hopeless in general. What changed recently is a framework for doing much better
    when you have \textbf{side information}.
  \end{itemize}
\end{frame}

\begin{frame}{The ``algorithms with predictions'' framework}
  \begin{block}{Mitzenmacher \& Vassilvitskii, \emph{Commun.\ ACM} 65(7), 2022, pp.\
  33--35. DOI: \texttt{10.1145/3528087}.}
    Augment a classical online algorithm with a (possibly wrong) machine-learned
    \textbf{prediction} about the future. Measure it on two axes:
    \begin{itemize}
      \item \textbf{Consistency}: the competitive ratio when the prediction is
      \emph{perfect}.
      \item \textbf{Robustness}: the competitive ratio in the \emph{worst case}, even
      if the prediction is completely wrong.
    \end{itemize}
  \end{block}
  \begin{itemize}
    \item The goal is a smooth trade-off: near-optimal (consistency close to $1$) when
    the forecast is good, but never worse than a fixed worst-case bound (robustness) if
    it is bad -- you never do worse than ``ignore the prediction and fall back to a
    classical online algorithm.''
  \end{itemize}
\end{frame}

\begin{frame}{Applied to exactly this week's problem}
  \small
  \begin{block}{Boyar, Favrholdt, Kamali \& Larsen -- ``Online Interval Scheduling with
  Predictions'' (WADS 2023; journal version arXiv:2302.13701)}
    Studies online interval scheduling where the algorithm is given a (possibly
    erroneous) prediction of the \emph{full request sequence}. They prove
    asymptotically \textbf{tight} upper and lower bounds on the achievable competitive
    ratio as a function of the prediction error, and tight consistency-vs-robustness
    trade-off curves -- then validate the theory on real scheduling workloads.
  \end{block}
  \begin{itemize}
    \item This is a direct descendant of today's lecture: same problem (maximize
    accepted non-overlapping intervals), same notion of ``compatible,'' but with the
    single-shot assumption of the offline greedy replaced by an online, predicted
    setting.
  \end{itemize}
\end{frame}

\begin{frame}{As current as it gets: a Nov.\ 2025 follow-up}
  \small
  \begin{block}{Antoniadis, Shahheidar, Shahkarami \& Soltani -- ``A Switching
  Framework for Online Interval Scheduling with Predictions'' (AAAI 2026; posted to
  arXiv Nov.\ 2025, arXiv:2511.16194)}
    Instead of committing to one predictor, maintain several candidate predictors
    simultaneously and \textbf{switch} between them adaptively as their track records
    diverge -- hedging against betting everything on a single, possibly bad, forecast.
  \end{block}
  \begin{itemize}
    \item This paper is still working its way through the AAAI 2026 publication
    pipeline as this deck is written -- about as close to ``the current edge of the
    field'' as a course supplement can responsibly cite.
    \item The core idea -- \emph{don't trust one prediction, hedge across several, and
    adapt} -- is a natural generalization of the exchange-argument mindset from today's
    lecture: keep whichever candidate is currently ``winning,'' and be willing to swap.
  \end{itemize}
\end{frame}

% =================================================================
\section{Summary}
% =================================================================

\begin{frame}{Summary}
  \footnotesize
  \begin{enumerate}
    \item \textbf{Matroids explain \emph{when} greedy works at all}: the Rado--Edmonds
    theorem shows exchange-argument-style correctness is not a coincidence for
    interval scheduling -- it is a consequence of an underlying matroid structure.
    \item \textbf{Submodularity is the natural next step} once ``maximize total
    weight'' becomes ``maximize with diminishing returns''; Kobayashi \& Terao (ICALP
    2024) show the classical $(1-1/e)$ quality bound can now be reached with far fewer
    oracle queries.
    \item \textbf{Removing the offline assumption reopens the problem}: with no
    information about the future, online interval scheduling has no constant
    competitive ratio; with a prediction, Boyar et al.\ (2023) give tight
    consistency/robustness trade-offs, and Antoniadis et al.\ (AAAI 2026) push this to
    hedging across multiple predictors.
    \item The throughline across all of it: \textbf{correctness first, then
    efficiency, then robustness to a changing model of the input} -- the same arc this
    course follows week over week.
  \end{enumerate}
  \begin{alertblock}{Looking ahead}
    Next week (Kleinberg \& Tardos, Ch.\ 4 continued) proves Dijkstra's algorithm and
    Kruskal/Prim's MST algorithms optimal with the \emph{same} two proof templates --
    and Union-Find from Week 3 reappears as Kruskal's core data structure.
  \end{alertblock}
\end{frame}

\end{document}
